\begin{proofbox}
\textbf{准备：统一记号与两个基础事实}

记 $r_1=PF_1$、$r_2=PF_2$，$F_1=(-c,0)$、$F_2=(c,0)$，
$\varepsilon=\operatorname{sgn}(x_0)$。
由于 $P$ 不在顶点，$|x_0|>a$ 且 $y_0\ne0$。
对双曲线上一点，有
\[
\boxed{r_1-r_2=2\varepsilon a},\qquad
\boxed{r_1+r_2=2e|x_0|}.
\tag{1}
\]
第一式是双曲线的定义；第二式可由焦半径公式得到，下面还会直接验证。

\textbf{方法一（首选最快）：内切圆切线长法——三行求 $x_I$}

设内切圆与 $F_1F_2$ 切于 $T$，与 $PF_1$、$PF_2$ 分别切于 $M,N$。
因同一点引圆的两条切线相等，
$F_1T=F_1M$、$F_2T=F_2N$、$PM=PN$。
把三条边长相加减，即
\[
2F_1T=F_1F_2+PF_1-PF_2=2c+(r_1-r_2).
\]
由于 $I$ 在 $T$ 的正上方或正下方，
\[
\boxed{x_I=x_T=-c+F_1T=\frac{r_1-r_2}{2}}
\overset{(1)}{=}\boxed{\varepsilon a}.
\]
\textbf{三步记法：同点切线等长 $\rightarrow$ 切点横坐标为半个焦距差 $\rightarrow$ 双曲线定义。}
这条证明不需要设出内心坐标，也不需要解二元方程。

\textbf{方法二（最快补出 $y_I$）：面积除以半周长}

将 $\triangle PF_1F_2$ 用三条连向内心的线段分成三个小三角形。
它们到原三角形三条边的高都等于内切圆半径 $r$，所以
\[
S=\frac12r(r_1+r_2+2c)=rs,\qquad s=c+e|x_0|.
\]
另一方面，由底边为 $2c$、高为 $|y_0|$，
\[
S=\frac12\cdot2c\cdot|y_0|=c|y_0|.
\]
故
\[
r=\frac{c|y_0|}{c+e|x_0|},\qquad
\boxed{y_I=\frac{cy_0}{c+e|x_0|}
=\frac{ay_0}{a+|x_0|}}.
\]
内心一定与 $P$ 处于 $x$ 轴同侧，所以 $y_I$ 与 $y_0$ 同号，不能直接把 $r$ 当作有符号的纵坐标。
将方法一、二合并，就得到完整的内心坐标。

\tcbline
\textbf{方法三（代数通法）：对边长加权的内心坐标公式}

任意三角形 $ABC$ 中，内心 $I$ 的位置向量为
\[
I=\frac{|BC|A+|CA|B+|AB|C}{|BC|+|CA|+|AB|}.
\]
理解方法：三个小三角形 $IBC,ICA,IAB$ 的高都等于内切圆半径，面积之比等于三个对边之比，故是上述加权平均。
令 $A=P$、$B=F_1$、$C=F_2$，得到
\[
I=\frac{2cP+r_2F_1+r_1F_2}{2c+r_1+r_2}.
\]
逐项取坐标，得到
\[
x_I=\frac{2cx_0+c(r_1-r_2)}{2c+r_1+r_2}
=\frac{c(x_0+\varepsilon a)}{c+e|x_0|}.
\]
注意 $x_0=\varepsilon|x_0|$、$e=c/a$，于是
\[
x_I=\frac{c\varepsilon(|x_0|+a)}{c(1+|x_0|/a)}
=\varepsilon a.
\]
纵坐标同理：
\[
y_I=\frac{2cy_0}{2c+r_1+r_2}
=\frac{cy_0}{c+e|x_0|}.
\]
\textbf{适用场景：} 已知三个顶点和三边长，要求精确内心坐标时，这是最稳的解析几何通法。

\textbf{方法四（角平分线视角）：先求角平分线交底边的点}

设 $P$ 的内角平分线与 $F_1F_2$ 相交于 $Q$。
角平分线定理给出
\[
\frac{F_1Q}{QF_2}=\frac{PF_1}{PF_2}=\frac{r_1}{r_2},\qquad
x_Q=\frac{cr_1-cr_2}{r_1+r_2}.
\]
代入式 (1)，得
\[
\boxed{Q\left(\frac{a^2}{x_0},0\right)}.
\]
因为内心 $I$ 在 $PQ$ 上，而方法一已经确定 $x_I=\varepsilon a$，
由同一直线上纵坐标的线性关系，
\[
\frac{y_I}{y_0}=
\frac{x_I-x_Q}{x_0-x_Q}
=\frac{\varepsilon a-a^2/x_0}{x_0-a^2/x_0}
=\frac{a}{a+|x_0|}.
\]
于是 $y_I=ay_0/(a+|x_0|)$。
\textbf{几何画图顺序：} 焦点连线 $\rightarrow$ 内角平分线 $PQ$
$\rightarrow$ 固定竖线 $x=\varepsilon a$ $\rightarrow$ 两线交点 $I$。
\textbf{注意：} $Q$ 是\emph{角平分线}与底边的交点，$T$ 是\emph{内切圆切点}，一般 $Q\ne T$！

\textbf{方法五（不直接套焦半径公式）：先从曲线方程推出 $r_1,r_2$}

由双曲线方程
\[
y_0^2=\frac{b^2}{a^2}(x_0^2-a^2),\qquad c^2=a^2+b^2,
\]
可得
\begin{align*}
r_1^2&=(x_0+c)^2+y_0^2
=\frac{c^2}{a^2}x_0^2+2cx_0+a^2
=(ex_0+a)^2,\\
r_2^2&=(x_0-c)^2+y_0^2
=\frac{c^2}{a^2}x_0^2-2cx_0+a^2
=(ex_0-a)^2.
\end{align*}
所以
\[
\boxed{r_1=|ex_0+a|},\qquad
\boxed{r_2=|ex_0-a|}.
\]
当 $x_0>a$ 时，$r_1=ex_0+a$、$r_2=ex_0-a$；
当 $x_0<-a$ 时，$r_1=-ex_0-a$、$r_2=-ex_0+a$。
这就严格推出式 (1)，同时解释了为什么\textbf{左支必须处理绝对值}。

\tcbline
\textbf{方法六（垂心最快）：只联立两条高线}

因为 $F_1F_2$ 在 $x$ 轴上，所以从 $P$ 作出的高线为
\[x=x_0.\]
直线 $PF_2$ 的斜率为 $y_0/(x_0-c)$，当 $x_0=c$ 时它是竖直线。
统一使用\textbf{垂直方向向量}，可直接写出过 $F_1$ 的高线：
\[
(x_0-c)(x+c)+y_0y=0.
\]
这里该方程对 $x_0=c$ 也成立，不必除以 $x_0-c$。
代入 $x=x_0$：
\[
(x_0-c)(x_0+c)+y_0y_H=0,
\]
即
\[
\boxed{H\left(x_0,\frac{c^2-x_0^2}{y_0}\right)}.
\]
\textbf{特别情形：} 当 $x_0=c$（右支）时，$H=F_2$；
当 $x_0=-c$（左支）时，$H=F_1$。
这是三角形在对应焦点处为直角三角形的自然结果。

\textbf{方法七（外心最快）：中垂线加一次平方相等}

由 $F_1F_2$ 的中点为原点且底边水平，可设外心为 $O'=(0,t)$。
外心到 $F_2$、$P$ 的距离相等：
\[
c^2+t^2=x_0^2+(y_0-t)^2.
\]
展开、约掉 $t^2$：
\[
2y_0t=x_0^2+y_0^2-c^2.
\]
故
\[
\boxed{O'\left(0,\frac{x_0^2+y_0^2-c^2}{2y_0}\right)}.
\]
\textbf{另一种思路：} 圆方程写为
$x^2+y^2+Dx+Ey+F=0$。
代入 $F_1,F_2$ 并相减，得 $D=0$；再代入 $P$，
得 $E=(c^2-x_0^2-y_0^2)/y_0$，
圆心 $( -D/2,-E/2)$ 就是上式。

\textbf{方法八（垂心第二证）：欧拉线一行推出}

重心直接取三点平均：
\[
\boxed{G=\frac{P+F_1+F_2}{3}=
\left(\frac{x_0}{3},\frac{y_0}{3}\right)}.
\]
任意非退化三角形的欧拉线关系是
\[
\boxed{H=3G-2O'}.
\]
其简单向量证明如下：以外心 $O'$ 为原点，令三个顶点的位置向量为
$\boldsymbol u,\boldsymbol v,\boldsymbol w$，长度都等于外接圆半径。
设 $\boldsymbol h=\boldsymbol u+\boldsymbol v+\boldsymbol w$，则
\[
(\boldsymbol h-\boldsymbol u)\cdot(\boldsymbol v-\boldsymbol w)
=(\boldsymbol v+\boldsymbol w)\cdot(\boldsymbol v-\boldsymbol w)
=|\boldsymbol v|^2-|\boldsymbol w|^2=0.
\]
即从顶点 $\boldsymbol u$ 指向 $\boldsymbol h$ 的直线垂直于对边；其余顶点也相同，
所以 $\boldsymbol h$ 就是垂心向量。
又重心向量 $\boldsymbol g=(\boldsymbol u+\boldsymbol v+\boldsymbol w)/3$，
因此 $H=3G-2O'$（还原到通常坐标后仍成立）。
现在代入已求得的 $G,O'$：
\[
3G-2O'=\left(x_0,y_0-\frac{x_0^2+y_0^2-c^2}{y_0}\right)
=\boxed{\left(x_0,\frac{c^2-x_0^2}{y_0}\right)}.
\]

\textbf{检验：} 内心必须在三角形内部；四心公式在左右支、$y_0$ 正负两种情况下都成立。
但 $y_0=0$ 时三角形退化，不能套用任何四心结论。
\end{proofbox}
